\documentclass[a4paper,12pt]{extarticle}
\usepackage{geometry}
\usepackage[T1]{fontenc}
\usepackage[utf8]{inputenc}
\usepackage[english,russian]{babel}
\usepackage{amsmath}
\usepackage{amsthm}
\usepackage{amssymb}
\usepackage{fancyhdr}
\usepackage{setspace}
\usepackage{graphicx}
\usepackage{colortbl}
\usepackage{tikz}
\usepackage{pgf}
\usepackage{subcaption}
\usepackage{listings}
\usepackage{indentfirst}
\usepackage[
backend=biber,
style=numeric,
maxbibnames=99
]{biblatex}
\addbibresource{refs.bib}
\usepackage[colorlinks,citecolor=blue,linkcolor=blue,bookmarks=false,hypertexnames=true, urlcolor=blue]{hyperref} 
\usepackage{indentfirst}
\usepackage{mathtools}
\usepackage{booktabs}
\usepackage[flushleft]{threeparttable}
\usepackage{tablefootnote}

\usepackage{chngcntr} % нумерация графиков и таблиц по секциям
\counterwithin{table}{section}
\counterwithin{figure}{section}

\graphicspath{{graphics/}}%путь к рисункам

\makeatletter
% \renewcommand{\@biblabel}[1]{#1.} % Заменяем библиографию с квадратных скобок на точку:
\makeatother

\geometry{left=2.5cm}% левое поле
\geometry{right=1.0cm}% правое поле
\geometry{top=2.0cm}% верхнее поле
\geometry{bottom=2.0cm}% нижнее поле
\setlength{\parindent}{1.25cm}
\renewcommand{\baselinestretch}{1.5} % междустрочный интервал


\newcommand{\bibref}[3]{\hyperlink{#1}{#2 (#3)}} % biblabel, authors, year
\addto\captionsrussian{\def\refname{Список литературы (или источников)}} 

\renewcommand{\theenumi}{\arabic{enumi}}% Меняем везде перечисления на цифра.цифра
\renewcommand{\labelenumi}{\arabic{enumi}}% Меняем везде перечисления на цифра.цифра
\renewcommand{\theenumii}{.\arabic{enumii}}% Меняем везде перечисления на цифра.цифра
\renewcommand{\labelenumii}{\arabic{enumi}.\arabic{enumii}.}% Меняем везде перечисления на цифра.цифра
\renewcommand{\theenumiii}{.\arabic{enumiii}}% Меняем везде перечисления на цифра.цифра
\renewcommand{\labelenumiii}{\arabic{enumi}.\arabic{enumii}.\arabic{enumiii}.}% Меняем везде перечисления на цифра.цифра

\begin{document}
\begin{titlepage}
	\newpage
	
	{\setstretch{1.0}
		\begin{center}
			ПРАВИТЕЛЬСТВО РОССИЙСКОЙ ФЕДЕРАЦИИ\\
			ФГАОУ ВО НАЦИОНАЛЬНЫЙ ИССЛЕДОВАТЕЛЬСКИЙ УНИВЕРСИТЕТ\\
			«ВЫСШАЯ ШКОЛА ЭКОНОМИКИ»
			\\
			\bigskip
			Факультет компьютерных наук\\
			Образовательная программа «Прикладная математика и информатика»
		\end{center}
	}
	
	\vspace{2em}
	\vspace{4em}
	
	\begin{center}
		{\bf Отчет о командном программном проекте на тему:}\\
		{\bf Разработка онлайн бизнес-игры “Мультикомандный Featureban”}\\
		% строчка ниже нужна только при сдача плана КР, при финальной сдаче закомментируйте ее
		% (промежуточный, этап 1)
	\end{center}
	
	\vspace{2em}
	
	{\bf Выполнили студенты: \vspace{2mm}}
	
	{\setstretch{1.1}
		\begin{tabular}{l@{\hskip 1.5cm}l}
			группы \#БПМИ244, 2 курса & Бакунц Арина Александровна \\
			группы \#БПМИ245, 2 курса & Денисов Николай Дмитриевич \\
			группы \#БПМИ246, 2 курса & Овчаров Иван Евгеньевич \\
			группы \#БПМИ249, 2 курса & Мостова Александра Александровна \\
			группы \#БПМИ249, 2 курса & Сухобок Тимофей Михайлович \\
			группы \#БПМИ2411, 2 курса & Чупрына Матвей Михайлович
	\end{tabular}}
	
	\vspace{1em}
	{\bf Принял руководитель проекта: \vspace{2mm}}
	
	{\setstretch{1.1}
		\begin{tabular}{l}
			Коротков Александр Олегович\\
			Руководитель группы, группа дополнительных продуктов\\
			АО "Т-Банк"
	\end{tabular}}
	
	\vspace{\fill}
	
	\begin{center}
		Москва 2026
	\end{center}
	
\end{titlepage}% это титульный лист - выберите подходящий вам из имеющихся в проекте вариантов (kr - курсовая работа у 3 курса, vkr - выпускная квалификационная работа у 4 курса)
\newpage
\setcounter{page}{2}

{
	\hypersetup{linkcolor=black}
	\tableofcontents
}

\newpage

\newpage
\section{Аннотация}   % this is how to use russian
Проект посвящен разработке веб-версии бизнес-игры «Мультикомандный Featureban» — инструмента для проведения дистанционных корпоративных тренингов, направленных на развитие навыков межкомандного взаимодействия и совершенствование коммуникации в условиях совместной работы над проектами.
Эта игра помогает просимулировать популярный в данный момент в крупных IT-компаниях метод разработки - Kanban~\cite{kanban-book-bor}.

Классическая очная версия игры давно зарекомендовала себя как эффективный способ выявления и разрешения проблем координации и синхронизации действий между разными подразделениями при разработке одного продукта. Однако её очный формат существенно ограничивает масштабируемость применения: требует физического присутствия участников (а также помещение для проведения тренинга), усложняет работу модератора и создает барьеры в виде географических ограничений.

Онлайн-версия устранит эти ограничения, позволяя проводить тренинги в полностью дистанционном формате, упростив процесс проведения данной игры и обеспечив возможность одновременного участия команд, расположенных в разных регионах. Приложение будет предоставлять игрокам интуитивный интерфейс для взаимодействия и отслеживания игровых метрик, а также специализированный модуль для тренера, позволяющий контролировать ход игровой сессии и анализировать результаты .

Планируемые результаты: Разработано полнофункциональное веб-приложение с архитектурой, обеспечивающей стабильное взаимодействие множества пользователей в реальном времени. Система включает интерфейс для участников (с визуализацией процессов, отображением карточек и метрик), модуль управления командами и специализированный инструмент для тренера. 


\section*{Ключевые слова}
Канбан, симуляция, бизнес-игра, Фичебан, мультикомандный, корпоративные тренинги, веб-приложение, закон Литтла, Телеграм бот
\pagebreak

\section{Введение}

В условиях современного развития информационных технологий и усложнения программных продуктов существенно возрастает сложность процессов их разработки и сопровождения. Современные компании, особенно в сфере информационных технологий, всё чаще переходят от работы отдельных автономных команд к модели, в которой несколько команд параллельно развивают один продукт или платформу. Такая организация труда требует более высокого уровня координации, прозрачности процессов и согласованности действий между участниками.

В этих условиях особое значение приобретают гибкие методологии управления проектами, ориентированные не только на локальную эффективность отдельных команд, но и на оптимизацию всей системы в целом. Одной из наиболее распространённых и практически применяемых методологий является Kanban. В отличие от итеративных подходов, таких как Scrum, Kanban ориентирован на непрерывный поток работ, визуализацию процессов и ограничение количества незавершённых задач. Данный подход хорошо масштабируется и широко используется в организациях со сложной структурой и уже существующими продуктами.

При этом ключевые принципы Kanban, включая ограничение WIP (Work in Progress), управление потоком и анализ метрик производительности, достаточно сложно усваиваются при изучении исключительно теоретических материалов. Без практического опыта участникам трудно осознать влияние управленческих решений на пропускную способность системы, время выполнения задач и стабильность процессов. В результате между формальным знанием методологии и её реальным применением на практике возникает существенный разрыв.

Для преодоления данного разрыва в корпоративном обучении и образовательных программах широко применяются бизнес-игры и игровые симуляции. Они позволяют участникам в безопасной среде экспериментировать с процессами, принимать управленческие решения и наблюдать их последствия без риска для реальных проектов. Игровой формат способствует более глубокому пониманию системных эффектов, возникающих при управлении потоком работ, и позволяет наглядно продемонстрировать влияние ограничений и изменений в процессе разработки.

Одной из таких симуляций является игра Featureban, моделирующая процесс разработки продукта с использованием Kanban-подхода. Классическая версия игры ориентирована на работу одной команды и позволяет продемонстрировать базовые принципы управления потоком. Мультикомандная версия Featureban представляет собой развитие данной симуляции и направлена на моделирование взаимодействия нескольких команд, одновременно работающих над общим продуктом. В рамках такой игры проявляются межкомандные зависимости, узкие места системы, проблемы синхронизации и эффекты локальной оптимизации, характерные для реальных крупных проектов.

Актуальность данной работы обусловлена возрастающей сложностью управления процессами разработки в условиях масштабирования продуктов и увеличения количества команд, вовлечённых в совместную работу. На практике многие организации сталкиваются с тем, что методы и практики Kanban, успешно применяемые на уровне одной команды, теряют эффективность при переходе к мультикомандному формату. При этом существующие инструменты и обучающие материалы не позволяют в полной мере продемонстрировать системные эффекты, возникающие при межкомандном взаимодействии, и обеспечить корректное воспроизведение игровой механики мультикомандной симуляции.

В этой связи актуальной является задача создания специализированного программного инструмента, позволяющего наглядно моделировать мультикомандные процессы разработки, демонстрировать влияние управленческих решений на всю систему в целом и обеспечивать автоматизацию игровых правил и сбора метрик. Решение данной задачи способствует повышению эффективности обучения и более осознанному внедрению Kanban-подходов в организациях со сложной структурой.

Целью данной курсовой работы является разработка онлайн-версии бизнес-игры «Мультикомандный Featureban», предназначенной для моделирования межкомандного взаимодействия и использования в обучающих и тренинговых форматах. Для достижения поставленной цели в рамках работы проводится анализ предметной области, формулируются требования к разрабатываемой системе, проектируется архитектура веб приложения и реализуются основные компоненты системы.

В результате планируется создать веб приложение, поддерживающее проведение игровых сессий с участием нескольких команд, разграничение ролей пользователей, визуализацию процессов и автоматизированный сбор ключевых метрик игрового процесса. Новизна данной работы заключается в разработке специализированного инструмента для онлайн-проведения мультикомандной версии Featureban, который учитывает особенности игровой механики, поддерживает сценарии симуляции и обеспечивает централизованное управление игровым процессом. Практическая значимость работы состоит в возможности использования разработанного приложения для проведения корпоративных тренингов и образовательных мероприятий, направленных на изучение принципов управления потоком и межкомандного взаимодействия.

\section{Обзор литературы}
Среди методов управления самыми популярными являются Scrum и Kanban. Идея методологии Scrum - разработка итерациями. То есть за определенный промежуток времени команда должна выпустить часть готового продукта. Scrum традиционно применяют при создании нового продукта. Kanban, напротив, чаще используют в крупных компаниях с уже работающим продуктом и множеством команд. Этот метод сфокусирован на управлении потоками работ и ограничении незавершенных работ.

Принципы гибких методов управления сложно усвоить исключительно с помощью книг и статей. Эта проблема подтверждается результатами исследований в области менеджмента и бизнес-образования. Например, обзор \textit{The effectiveness of games for educational purposes: A review of recent research}~\cite{games-for-edu} подтверждает: игровые методы и симуляции развивают управленческие навыки лучше, чем лекции или чтение литературы. Это связано с необходимостью понимать внутренние процессы системы и их взаимосвязь.

В первую очередь это касается методологий Kanban~\cite{kanban-book-bor} и Lean~\cite{lean-book}. Поэтому в обучении широко применяются игровые симуляции. В игре участники могут тестировать любые решения и сразу видеть их последствия, не боясь навредить реальному проекту. Эффективность этого метода подтверждается теорией из книги Д. Колба \textit{Experiential Learning Cycle}~\cite{kolb}. Согласно его модели, глубокое обучение проходит через четыре этапа: накопление конкретного опыта, рефлексивное наблюдение, абстрактную концептуализацию и активное экспериментирование. Бизнес-игра позволяет быстро пройти этот цикл, закрепляя ценные практические навыки.

В работе \textit{Learning with a strategic management simulation game: A case study}~\cite{learning-with-simulation} авторы также изучали эффективность бизнес-симуляций. Исследователи делают вывод, что игры помогают увидеть картину целиком и использовать теорию в деле. Так проще понять связи внутри процесса, чем при изучении обычных учебных материалов.

Рассмотрим подробнее эти подходы. В своей книге Дэвид Андерсон определяет шесть базовых практик метода Kanban~\cite{kanban-book-and}:
\begin{enumerate}
    \item Визуализация потока задач.
    \item Ограничение объема незавершенной работы (WIP).
    \item Непосредственное управление потоком.
    \item Четкие правила выполнения работы.
    \item Регулярные встречи для сбора обратной связи.
    \item Совместное улучшение процесса.
\end{enumerate}
а также подчеркивает принцип: \textit{start with what you do now}. В игровой форме участникам наглядно видно связь принятых решений в управлении и скорость доведения задач до конца. Этот метод тесно связан с концепцией Lean. Целью обоих методов является устранение потерь и увеличение пропускной способности. Практическая часть опирается на Теорию ограничений Э. Голдратта~\cite{gold-goal-book}. Согласно ей, общая производительность системы лимитируется её "узкими местами". Эти принципы легли в основу разрабатываемой симуляции, моделирующей реальные процессы для обучения. 

Также в управлении проектами часто используют закон Литтла~\cite{little-law}. На практике этот закон позволяет менеджерам проектов точнее прогнозировать сроки и понимать реальную скорость работы команды. В методологии Kanban закон выражается формулой:
\begin{equation}
    \label{little-law-formula}
    L = \lambda\times W
\end{equation}
где $L$ - это незавершенная работа (количество активных задач в системе), $\lambda$ - пропускная способность (количество выполненных задач за единицу времени), $W$ - время цикла (время прохождения задачи через процесс). Однако, прочитав, сложно понять, как она работает на самом деле, людьми эта формула воспринимается абстрактно. Игровая симуляция дает людям понять на примере, как действует этот закон, и позволяет увидеть, какой результат дают ограничения на WIP. 

На данный момент есть онлайн-версия игры Featureban\footnote{Онлайн версия игры доступна по ссылке: \url{https://tools.flowfast.io/featureban}, дата обр. 13.01.2026}, но эта версия игры только для одной команды и рассчитана на 3–5 человек. Версия игры на несколько команд существует только в оффлайн формате и разработана сертифицированными экспертами из России (Neogenda)\cite{multi-featureban}. где $L$ - это незавершенная работа (количество активных задач в системе), $\lambda$ - пропускная способность (количество выполненных задач за единицу времени), $W$ - время цикла (время прохождения задачи через процесс). Однако, прочитав, сложно понять, как она работает на самом деле, людьми эта формула воспринимается абстрактно. Игровая симуляция дает людям понять на примере, как действует этот закон, и позволяет увидеть, какой результат дают ограничения на WIP. 

На данный момент есть онлайн-версия игры Featureban\footnote{Онлайн версия игры доступна по ссылке: \url{https://tools.flowfast.io/featureban}, дата обр. 13.01.2026}, но эта версия игры только для одной команды и рассчитана на 3–5 человек. Версия игры на несколько команд существует только в оффлайн формате и разработана сертифицированными экспертами из России (Neogenda)\cite{multi-featureban}. Исходная версия Featureban от Майка Барроуза доступна под лицензией Creative Commons Attribution ShareAlike 4.0~\cite{license-cc}. Все материалы разрабатываемой мультикомандной модификации наследуют этот же правовой статус. Сейчас для проведения игры с несколькими командами приходится адаптировать сторонние платформы, например, Miro или Trello. Однако данные решения не могут быть напрямую использованы для реализации симуляции. 

Во-первых, они только визуализируют игру, но не накладывают никаких ограничений на действия участников. Из-за этого пропадает поддержка ролей, разделение на команды, разрушаются четкие сценарии игр, что являются критически важными аспектами симуляции для корректного моделирования принципов Kanban. Исследования игровых симуляций подчеркивают, что строгость правил, ограничений и сценариев является критически важным аспектом подобных обучающих методов. 

Во-вторых, подобные платформы не собирают нужную статистику. При отсутствии автоматизации сбора статистики увеличивается влияние человеческого фактора, становится сложнее проводить рефлексию над экспериментом, что приводит к снижению педагогической ценности проведения этой игры. 

Таким образом, такие универсальные инструменты не предоставляют достаточно функциональный интерфейс для высокоэффективного проведения мультикомандной версии игры.

\newpage 
\section{План дальнейшей работы}


\subsection{1. Этапы реализации проекта}


В рамках курсового проекта планируется последовательная реализация следующих этапов:

\begin{enumerate}
\item Анализ предметной области и формирование требований к системе.
\item Проектирование архитектуры онлайн-версии бизнес-игры.
\item Реализация серверной части системы.
\item Реализация клиентской части системы.
\item Интеграция компонентов и тестирование.
\item Подготовка демонстрационного сценария и сбор обратной связи.
\end{enumerate}

Указанные этапы могут уточняться и дополняться по мере развития проекта.

\subsection{2. Функциональная декомпозиция системы}


Проектируемая система включает следующие основные функциональные составляющие:

\begin{itemize}
\item \textbf{Игровая логика}

\begin{itemize}
\item сценарий проведения игры;
\item этапы игры и переходы между ними;
\item сбор и расчёт игровых метрик (WIP, lead time и др.).
\end{itemize}
\end{itemize}

\begin{itemize}
\item \textbf{Пользовательские роли}

\begin{itemize}
\item ведущий (модерация игры, управление процессом);
\item команда (совместное принятие решений);
\item игрок (участие в игровых активностях).
\end{itemize}
\end{itemize}

\begin{itemize}
\item \textbf{Сетевая и многокомандная часть}

\begin{itemize}
\item синхронизация состояния игры между участниками;
\item поддержка нескольких команд в рамках одной сессии.
\end{itemize}
\end{itemize}

\begin{itemize}
\item \textbf{Пользовательский интерфейс}

\begin{itemize}
\item интерфейс ведущего;
\item интерфейс команд и игроков.
\end{itemize}
\end{itemize}

\begin{itemize}
\item \textbf{Хранение и обработка данных}

\begin{itemize}
\item данные игровых сессий;
\item результаты игр и статистика.
\end{itemize}
\end{itemize}

\subsection{3. Предварительное распределение задач по участникам}


Распределение задач между участниками проекта носит предварительный характер и может изменяться в ходе выполнения работы.

\begin{itemize}
\item Анализ предметной области, формирование требований.
\item Проектирование архитектуры системы.
\item Реализация серверной логики и игровых механик.
\item Реализация клиентской части и пользовательских интерфейсов.
\item Интеграция компонентов и тестирование.
\end{itemize}

Конкретное распределение ответственности уточняется по мере развития проекта.


\newpage 
\printbibliography[heading=bibintoc] 

% \begin{thebibliography}{0}
% 	\bibitem{chirkova18}\hypertarget{chirkova18}{}
% 	\href{https://arxiv.org/abs/1810.10927}
% 	{Nadezhda Chirkova, Ekaterina Lobacheva, Dmitry Vetrov. Bayesian Compression for Natural Language Processing. In EMNLP 2018.}
% \end{thebibliography}

\end{document}
